\scanpage{75}
\deleted{\unsure{Écrire}}

Or dans le cas qui nous occupe \inserted{(si $\omega\in L$ est une $1$-forme invariante)}, on aura
\[
\begin{aligned}
 \beta(\omega,\omega)
 &=\int_X\omega\bar\omega
 =\int_X\underset{\substack{x=\unsure{\mathrm R|\Re|\mathrm{Re}}z\\y=\mathrm{Im}\,z}}{(d\unsure{x}+i\,dy)(dx-i\,dy)}
 =-2i\int_Xdx\wedge dy=-2i\lvert\omega_1\rvert^2\mathrm{Im}\,\tau\\
 &\hphantom{{}=\int_X\omega\bar\omega\qquad\qquad\qquad\qquad\qquad}=\lvert\omega_1\rvert^2(\bar\tau-\tau)
\end{aligned}
\]
La condition \unsure{résulte|visuelle|équivaut|voulue} \unsure{déjà|donc} au fait que
\[
 \beta=\deleted{\unsure{\rho}}-2i\rho\qquad\text{avec}\quad\rho>0\qquad \rho=\int_X dx\wedge dy
\]

\unsure{Quand|Quant} à $\alpha$, il reste \unsure{à|a} priori un \unsure{élém\textsuperscript{t}} arbitraire de $\overline L^{-1}\otimes L$. Notons que ``$\alpha$'' est défini, \inserted{(\unsure{indépendament|indépendamment} du choix d'une base $\omega$ de $L=\underline\omega$)} qui est-il ?

[N.B. \quad $\gamma=-\alpha\beta^{-1}\in L^{\otimes2}$ \unsure{.} est l'objet qui fait passer du splitting ``algébrique'' de $V$ au ``splitting transcendant''.

La connaissance de $(\alpha,\beta)$ (i.e. de $x\unsure{\mathrel{\sim}|\mapsto}\bar x$ \inserted{dans $V$}) équivaut à celle de $\underline\beta$ et $\underline\gamma$, i.e. de la \inserted{forme-}volume $\int_X\omega\bar\omega$ \uline{et} du splittage. ]

Posons, si $\omega,y$ sont choisis [donc $L\xrightarrow{\sim}\mathbb C$, $L^{-1}\unsure{\hookrightarrow|\xrightarrow{\sim}|\simeq}\mathbb C$ isomorphismes donnés), en plus de $\deleted{\rho=\mathrm{Vol}(X)}$,
\[
\begin{aligned}
 &\rho=-\frac1{2i}\int\omega\bar\omega=\mathrm{Vol}(X),\\
 \text{la fonct\textsuperscript{n}}\quad &\sigma=-\frac1{2i}\int y_0\bar y_0
\end{aligned}
\]
où $y_0$ est une forme différentielle fermée dont la classe est celle de $y$

On trouve \qquad $4\rho\sigma=(1-\alpha\bar\alpha)$
% archive sheet number: 74
